% !TEX program = pdflatex
\documentclass[11pt,a4paper]{article}
\usepackage[utf8]{inputenc}
\usepackage[T1]{fontenc}
\usepackage[spanish,es-nodecimaldot]{babel}
\usepackage[a4paper,margin=2.35cm]{geometry}
\usepackage{fancyhdr}
\usepackage{lmodern,xcolor,amsmath,amssymb,mathtools,tcolorbox}
\usepackage[hidelinks]{hyperref}
\usepackage{booktabs,array,tikz}
\usetikzlibrary{decorations.pathreplacing}
\usepackage{enumitem}
\setlist[itemize]{topsep=0pt, partopsep=0pt}
\definecolor{inkred}{HTML}{E84C4F}
\definecolor{palered}{HTML}{FFF1F1}
\definecolor{pencilgray}{HTML}{999999}
\definecolor{markpink}{HTML}{F2A7AC}
\definecolor{markyellow}{HTML}{F3D66E}
\pagestyle{plain}
\makeatletter
\renewcommand\section{\@startsection{section}{1}{\z@}{-1em}{0.4em}
 {\normalfont\large\bfseries\color{inkred}}}
\renewcommand\subsection{\@startsection{subsection}{2}{\z@}{-0.7em}{0.25em}
 {\normalfont\normalsize\bfseries\color{inkred}}}
\makeatother
\setlength{\parindent}{0pt}\setlength{\parskip}{0.5\baselineskip}
\setlength{\leftmargini}{1.35em}
\newenvironment{apuntebox}[1]{\begin{tcolorbox}[colback=palered,colframe=inkred,
 boxrule=0.7pt,arc=0mm,left=0.8em,right=0.8em,top=0.45em,bottom=0.45em]
 \textcolor{inkred}{\textbf{#1}}\quad}{\end{tcolorbox}}
\newcommand{\rojo}[1]{\textcolor{inkred}{#1}}
\newcommand{\gris}[1]{\textcolor{pencilgray}{#1}}
\newcommand{\E}{\mathbb{E}}\newcommand{\Prob}{\mathbb{P}}
\newcommand{\argmax}{\operatorname*{arg\,max}}
\newcommand{\Th}{\Theta}
\tikzset{decision/.style={circle,fill=black,inner sep=1.5pt},
 nature/.style={circle,draw=black,inner sep=1.5pt},
 branch/.style={draw=black,line width=.55pt},
 infoset/.style={draw=black,dashed,line width=.65pt},
 lab/.style={font=\small},paylab/.style={font=\scriptsize,align=center}}

%----------------------------------------------------------------------------------------
% HEADER
%----------------------------------------------------------------------------------------

\pagestyle{fancy}
\fancyhf{}

\fancyhead[L]{Octavio Sivori}
\fancyhead[R]{Micro II}
\fancyfoot[C]{\thepage}

\renewcommand{\headrulewidth}{0.4pt}
\renewcommand{\footrulewidth}{0pt}
\setlength{\headheight}{14.5pt}

%----------------------------------------------------------------------------------------------------------------------------------------------------------------------------------%

\usepackage{graphicx}
\AtBeginDocument{\shorthandoff{<>}}
\usetikzlibrary{svg.path}
\definecolor{notesgold}{HTML}{BDA126}
\definecolor{notesorange}{HTML}{E99820}
\newcommand{\gold}[1]{\textcolor{notesgold}{#1}}
\newcommand{\orange}[1]{\textcolor{notesorange}{#1}}
\newcommand{\fuente}[1]{\par{\small\gris{#1}}\par}
\newenvironment{notebox}{\begin{tcolorbox}[colback=white,colframe=notesgold,boxrule=0pt,leftrule=.9pt,arc=0pt,left=.7em,right=.3em,top=.3em,bottom=.3em]}{\end{tcolorbox}}
\newcommand{\best}[1]{\begingroup\setbox0=\hbox{$#1$}\textcolor{red}{\underline{\textcolor{black}{\box0}}}\endgroup}
\newcommand{\cellcirc}[1]{\tikz[baseline=(p.base)]{\node[draw=red,ellipse,inner xsep=3pt,inner ysep=1pt](p){$#1$};}}
\usetikzlibrary{shapes.geometric}
\begin{document}
\begin{center}
{\LARGE\bfseries\rojo{Juegos Repetidos}}\quad
\gris{($\subset$ Juegos dinámicos)}
\end{center}
\fuente{Apuntes manuscritos: 5. Juegos IV, páginas 1--5.}
La principal pregunta es si se puede sostener como equilibrio algo distinto al EN estático.

\section{Juegos repetidos finitas veces}
En la última etapa, tienen que jugar un EN sí o sí. No hay incentivos a ``seguir cooperando''.
\begin{apuntebox}{Proposición.} Si un juego $G$ es tal que \rojo{$\exists!$ EN} y se repite un número \rojo{finito} de veces, el \rojo{único ESP} es el EN de $G$.
\end{apuntebox}
¿Qué pasa cuando el EN \underline{no es único}?

Supongamos:
\begin{center}\renewcommand{\arraystretch}{1.35}
$\begin{array}{c|c|c|c|}
 &L_2&M_2&R_2\\\cline{2-4}
L_1&\cellcirc{1,1}&\best{5},0&0,0\\\cline{2-4}
M_1&0,\best{5}&4,4&0,0\\\cline{2-4}
R_1&0,0&0,0&\cellcirc{3,3}\\\cline{2-4}
\end{array}$
\end{center}
Y quisiéramos sostener $(M_1,M_2)$ dado que $\max\sum_i V_i$, $\forall i$.

Podemos proponer:
\begin{quote}
$s_1$: empezar con $M_1$. Si en la primera etapa el resultado fue $(M_1,M_2)$, juega $R_1$. Si no, $L_1$.

$s_2$: empezar con $M_2$. Si en la primera etapa el resultado fue $(M_1,M_2)$, juega $R_2$. Si no, $L_2$.
\end{quote}
Y garantizamos que en la última etapa juegan un EN. Hace falta chequear que es un EN del juego completo.
\begin{center}\renewcommand{\arraystretch}{1.3}
$\begin{array}{c|c|c|c|}
 &L_2&M_2&R_2\\\cline{2-4}
L_1&1+1,1+1&5+1,0+1&0+1,0+1\\\cline{2-4}
M_1&0+1,5+1&4+3,4+3&0+1,0+1\\\cline{2-4}
R_1&0+1,0+1&0+1,0+1&3+1,3+1\\\cline{2-4}
\end{array}$
\par$\Downarrow$\par
$\begin{array}{c|c|c|c|}
 &L_2&M_2&R_2\\\cline{2-4}
L_1&\best{2},\best{2}&6,1&1,1\\\cline{2-4}
M_1&1,6&\cellcirc{7,7}&1,1\\\cline{2-4}
R_1&1,1&1,1&\best{4},\best{4}\\\cline{2-4}
\end{array}$
\end{center}
$\Rightarrow$ EN: $\{(M_1,M_2),\ldots\}$. Podemos sostener en alguna etapa algo $\neq$ al EN estático.

\section{Juegos repetidos infinitas veces}
Arrancamos con el dilema del prisionero donde los jugadores descuentan la utilidad con $\delta$.
\begin{center}\renewcommand{\arraystretch}{1.3}
\begin{tabular}{c c|c|c|}
&&$c$&$-c$\\\cline{3-4}
$(1)$&$c$&$\cellcirc{-4,-4}$&$\best{0},-5$\\\cline{3-4}
&$-c$&$-5,\best{0}$&$-1,-1$\\\cline{3-4}
\end{tabular}\quad $(2)$
\end{center}
Si juegan siempre $(c,c)$:
\[
S_n=-4+(-4)\delta+(-4)\delta^2+\cdots=\frac{-4}{1-\delta}.
\]
Si juegan siempre $(-c,-c)$:
\[
S_n=-1+(-1)\delta+(-1)\delta^2+\cdots=\frac{-1}{1-\delta}.
\]
¿Podemos conseguir que jueguen siempre $(-c,-c)$?

\subsection{Estrategias trigger}
\begin{apuntebox}{Definición. Estrategia trigger} En un juego infinito, una \rojo{estrategia trigger} para el jugador $i$ es:
\[
s_i(h^t)=
\begin{cases}
a_i^c &\text{si nadie se desvió en }T=1,\ldots,t-1,\\
a_i^p &\text{si algún desvío fue observado en algún }T.
\end{cases}
\]
\gris{$h^t$: historia en $t$.}
\end{apuntebox}
Estrategia trigger en el dilema del prisionero:
\begin{quote}
Empieza con $-c$. Si en todos los períodos anteriores el resultado fue $(-c,-c)$, juega $c$. Si en algún período fue distinto, juega $c$.
\end{quote}
Chequeamos que $\nexists$ incentivos al desvío con la $s_i$:
\[
u_i^{eq}=\frac{-1}{1-\delta},\qquad
u_i^{\text{desvío}}=0-4\delta+(-4)\delta^2+\cdots=\frac{-4\delta}{1-\delta}.
\]
No se desvía si
\[
\frac{-1}{1-\delta}\geq\frac{-4\delta}{1-\delta}
\quad\Longrightarrow\quad 1\leq4\delta
\quad\Longrightarrow\quad\delta\geq\frac14.
\]
Si $\delta\geq\frac14$, la estrategia trigger constituye un EN.

\underline{Otro ejemplo:} Cournot repetido infinitas veces.
\begin{quote}
Empiezo con $\dfrac{q^M}{2}=\dfrac{a-c}{4}$.

Si en todas las etapas el resultado fue $\left(\dfrac{a-c}{4},\dfrac{a-c}{4}\right)$, sigo produciendo $\dfrac{a-c}{4}$.

Si no, produzco $\dfrac{a-c}{3}$.
\end{quote}
Constituye un EN si $\delta\geq9/17$.

\underline{Otro ejemplo:} Bertrand repetido infinitas veces.
\begin{quote}
Empiezo con $p_i=\dfrac{a+c}{2}$.

Si en todas las etapas el resultado fue $\left(\dfrac{a+c}{2},\dfrac{a+c}{2}\right)$, sigo jugando $\dfrac{a+c}{2}$.

Si no, juego $c$.
\end{quote}
Constituye un EN si $\delta\geq1/2$.

\subsection{Pagos socialmente posibles y teoremas de Folk}
\begin{apuntebox}{Teorema de Folk} Si $G$ es un juego \rojo{estático, finito} y con \rojo{información completa}, cualquier vector de pagos \rojo{socialmente posible} donde cada jugador gana más que lo que gana en un EN estático puede ser sostenido como un \rojo{ESP} del juego repetido \rojo{infinitas veces} si $\delta\simeq1$.
\end{apuntebox}
\begin{apuntebox}{Definición. Pago socialmente posible} Un pago socialmente posible es una combinación convexa de todos los pagos del juego.
\end{apuntebox}
En el dilema del prisionero, los pagos socialmente posibles donde se gana más que en el EN están dados por:
\begin{center}
\begin{tikzpicture}[x=.9cm,y=.9cm,font=\small]
\draw[->](-5.8,0)--(.85,0)node[right]{$u_1$};
\draw[->](0,-5.8)--(0,.8)node[above]{$u_2$};
\fill[notesgold,opacity=.28](-4,-.25)--(-1,-1)--(-.25,-4)--(-4,-4)--cycle;
\draw[notesgold,line width=1.2pt](-4,-.25)--(-4,-4)--(0,-4);
\draw[red,line width=1pt](-5,0)--(-1,-1)--(0,-5)--(-4,-4)--cycle;
\foreach \x in {-4,-1}{\draw(\x,.1)--(\x,-.1)node[above=6pt]{$\x$};}
\foreach \y in {-4,-1}{\draw(-.1,\y)--(.1,\y)node[right=3pt]{$\y$};}
\fill(-5,0)circle(1.8pt)node[above left]{$(-5,0)$};
\fill(-1,-1)circle(1.8pt)node[below left=3pt]{$(-1,-1)$};
\fill(-4,-4)circle(1.8pt)node[below left]{$(-4,-4)$};
\fill(0,-5)circle(1.8pt)node[right]{$(0,-5)$};
\end{tikzpicture}
\end{center}
\begin{apuntebox}{Teorema de Fudenberg y Maskin} Si la dimensión del conjunto de pagos socialmente posibles es igual al número de jugadores, entonces \rojo{cualquier pago} socialmente posible en el que \rojo{todos} los jugadores reciban más que su \rojo{MINMAX} puede sostenerse como un \rojo{ESP} del juego repetido $\infty$ veces si son suficientemente pacientes.
\end{apuntebox}
\subsection{MINMAX}
\begin{apuntebox}{Definición. MINMAX} MINMAX es la utilidad de reserva.
\[
\min_{\sigma_{-i}}\max_{\sigma_i} V_i(\sigma_i,\sigma_{-i})
=\min_{\sigma_{-i}}V_i\bigl(\sigma_i(\sigma_{-i}),\sigma_{-i}\bigr).
\]
\end{apuntebox}
\gris{Ejemplo:}
\begin{center}\renewcommand{\arraystretch}{1.25}
$\begin{array}{c|c|c|}&A&B\\\cline{2-3}1&-2&1\\\cline{2-3}2&1&-2\\\cline{2-3}3&0&0\\\cline{2-3}\end{array}$
\end{center}
Supongamos que $J_2$ castiga al $J_1$, no le importan sus pagos. El MINMAX para 1 surge de:
\[
v_1=\min_{\sigma_2}\max_{\sigma_1}u_1(\sigma_1,\sigma_2).
\]
I)\quad $u_1(1\mid q)=-2q+1(1-q)=1-3q$,
\[
u_1(2\mid q)=1q-2(1-q)=3q-2,\qquad u_1(3\mid q)=0.
\]
\[
\Rightarrow\quad BR_1(q)=\max\{1-3q,3q-2,0\}.
\]
II)\quad $J_2$ resuelve
\[
\min_{\{q\}}\max\{1-3q,3q-2,0\}.
\]
\[
\Longleftrightarrow\quad q\in\left[\frac13,\frac23\right].
\]
\begin{center}
\begin{tikzpicture}[x=5cm,y=2.5cm,font=\small]
\draw[->](0,0)--(1.15,0)node[right]{$q$};
\draw[->](0,0)--(0,1.4)node[above]{$u_1$};
\draw[line width=.8pt](0,1)--(1/3,0)--(2/3,0)--(1,1);
\draw(-.025,1)--(.025,1)node[left=5pt]{$1$};
\foreach \x/\l in {0.333333/{1/3},0.666667/{2/3}}{\draw(\x,.03)--(\x,-.03)node[below=3pt]{$\l$};}
\end{tikzpicture}
\end{center}
\underline{Comments:}
\begin{itemize}
\item MINMAX es el máximo castigo que \#2 le puede imponer a \#1.
\item En \rojo{TODO EN estático}, cada jugador recibe al menos su MINMAX.
\item En el \rojo{juego repetido $\infty$ veces}, si el jugador juega la BR del estático, no puede recibir menos que el MINMAX.
\item En \rojo{ESP} NADIE recibe menos que el MINMAX.
\end{itemize}
\end{document}
